\chapter{9-10-26}

\section{Ecuaciones diofánticas: solución general}

\begin{proposicion}
Sean $a,b,c\in\mathbb{Z}$ con $a$ y $b$ no ambos nulos, y sea $d=(a:b)$. La
ecuación
\[
ax+by=c
\]
tiene solución entera si y sólo si $d\mid c$.

En ese caso, si $u,v\in\mathbb{Z}$ cumplen $ua+vb=d$ (o, equivalentemente,
$u\frac{a}{d}+v\frac{b}{d}=1$), entonces
\[
(x_0,y_0)=\left(\frac{c}{d}\,u,\ \frac{c}{d}\,v\right)
\]
es una solución particular, y el conjunto de todas las soluciones es
\[
S=\left\{\left(x_0-k\frac{b}{d},\ y_0+k\frac{a}{d}\right) : k\in\mathbb{Z}\right\}.
\]
\end{proposicion}

\begin{nota}{.}
Si $a=b=0$ la ecuación es $0=c$: no tiene solución si $c\neq 0$ y todo
$(x,y)\in\mathbb{Z}^2$ es solución si $c=0$. Por eso pedimos que algún
coeficiente sea no nulo.
\end{nota}

\begin{demostracion}
($\Rightarrow$) Como $d\mid a$ y $d\mid b$, si $ax+by=c$ entonces $d\mid c$.

($\Leftarrow$) Si $c=dc'$, y $ua+vb=d$ (Bézout), entonces
$a(uc')+b(vc')=dc'=c$, así que $(x_0,y_0)=(uc',vc')$ es solución particular.

\textbf{Todas las soluciones.} Sean $a'=a/d$ y $b'=b/d$, que cumplen $(a':b')=1$.
Cada elemento de $S$ es solución, pues
\[
a\left(x_0-k\tfrac{b}{d}\right)+b\left(y_0+k\tfrac{a}{d}\right)=ax_0+by_0-kab'+kba'=c,
\]
ya que $ab'=a\frac{b}{d}=ba'$. Recíprocamente, sea $(x,y)$ una solución
(suponemos $b\neq 0$; si $b=0$ se intercambian los roles de $x$ e $y$).
Restando $ax_0+by_0=c$:
\[
a(x-x_0)+b(y-y_0)=0 \iff a'(x-x_0)=-b'(y-y_0).
\]
Entonces $b'\mid a'(x-x_0)$ y, como $(a':b')=1$, por el lema de Gauss
$b'\mid x-x_0$, o sea $x-x_0=-kb'$ para algún $k\in\mathbb{Z}$. Reemplazando,
$a'(-kb')=-b'(y-y_0)$ y, cancelando $b'\neq 0$, $y-y_0=ka'$. Luego
$(x,y)\in S$.
\end{demostracion}

\begin{ejemplo}
Resolver $23x+11y=7$. Tenemos $d=(23:11)=1\mid 7$. Como
$23\cdot 1+11\cdot(-2)=1$, es $u=1$, $v=-2$ y
\[
(x_0,y_0)=(7\cdot 1,\ 7\cdot(-2))=(7,-14).
\]
Entonces $S=\{(7-11k,\ -14+23k): k\in\mathbb{Z}\}$, que es lo que obtuvimos en el
apunte del 02-10-26.
\end{ejemplo}

\section{La ecuación de congruencia $ax\equiv 0 \pmod m$}

\begin{proposicion}
Sean $a\in\mathbb{Z}$, $m\in\mathbb{Z}$ con $m\ge 1$, $d=(a:m)$ y $m'=m/d$.
Entonces el conjunto de todas las soluciones de $ax\equiv 0 \pmod m$ es
\[
\{km' : k\in\mathbb{Z}\}.
\]
Siempre hay soluciones (por ejemplo $x=0$).
\end{proposicion}

\begin{demostracion}
Sean $a'=a/d$ y $m'=m/d$, con $(a':m')=1$. Entonces
\[
m\mid ax \iff dm'\mid da'x \iff m'\mid a'x \iff m'\mid x,
\]
donde la segunda equivalencia es cancelar $d>0$ y la tercera usa
$(a':m')=1$ y el lema de Gauss. Finalmente $m'\mid x\iff x=km'$ con
$k\in\mathbb{Z}$.
\end{demostracion}

\begin{nota}{}
Si $m\mid a$, entonces $d=m$, $m'=1$ y todo $x\in\mathbb{Z}$ es solución: es
claro, pues $m\mid a\mid ax$.
\end{nota}

\begin{ejemplo}
Resolver $6x\equiv 0 \pmod{15}$:
\[
15\mid 6x \iff 15\mid 3\cdot 2x \iff 5\mid 2x \iff 5\mid x.
\]
Con la fórmula: $d=(6:15)=3$, $m'=15/3=5$, y las soluciones son $x=5k$.
\end{ejemplo}

\begin{nota}{}
Si $x$ es solución e $y\equiv x \pmod m$, entonces $y$ también es solución, pues
$ay\equiv ax\equiv 0 \pmod m$. Así que el conjunto de soluciones es una unión
de clases de congruencia módulo $m$.
\end{nota}

\begin{proposicion}
Con la notación anterior, las soluciones de $ax\equiv 0 \pmod m$ forman
exactamente $d$ clases módulo $m$, a saber
\[
[0],\ [m'],\ [2m'],\ \dots,\ [(d-1)m'].
\]
En particular, si $(a:m)=1$ la única clase solución es $[0]$, es decir
$x\equiv 0 \pmod m$.
\end{proposicion}

\begin{demostracion}
Sea $x=km'$ una solución, y dividimos $k$ por $d$: $k=qd+s$ con $0\le s<d$.
Entonces
\[
x=km'=qdm'+sm'=qm+sm',
\]
así que $x\equiv sm' \pmod m$, y $0\le sm'\le (d-1)m'=m-m'<m$. Luego toda solución
está en alguna de las clases $[sm']$ con $0\le s\le d-1$. Estas clases son
distintas entre sí, porque los números $0,m',\dots,(d-1)m'$ son distintos y
están entre $0$ y $m-1$ (restos distintos módulo $m$).
\end{demostracion}

\begin{ejemplo}
\begin{itemize}
    \item $3x\equiv 0 \pmod{15}$: $d=3$, $m'=5$. Soluciones: $x=5k$. Clases
    módulo $15$: $[0],[5],[10]$ (son $d=3$).
    \item $5x\equiv 0 \pmod{15}$: $d=5$, $m'=3$. Soluciones: $x=3k$. Clases
    módulo $15$: $[0],[3],[6],[9],[12]$ (son $d=5$).
    \item $2x\equiv 0 \pmod{16}$: $d=2$, $m'=8$. Soluciones: $x=8k$. Clases
    módulo $16$: $[0],[8]$ (son $d=2$).
\end{itemize}
\end{ejemplo}